\clearpage
\VerdierPhysicalPageAnchor{0236}%
\begingroup\small\setlength{\parskip}{0pt}\setlength{\emergencystretch}{1.5em}
\setlength{\abovedisplayskip}{4pt plus 2pt minus 2pt}%
\setlength{\belowdisplayskip}{4pt plus 2pt minus 2pt}%

\VerdierSemanticAnchor{ch4d-4-2-4}
\noindent\textbf{4.2.4.}\quad
Un morphisme d'objets filtrés de type $\widetilde{\mathbb Z}$ de
$\operatorname{comp}(\mathcal A)$ induit de la manière évidente un
morphisme entre les objets spectraux associés décrits par la proposition
(4.2.3), a). La proposition (4.2.3) définit donc un foncteur de la catégorie
des objets filtrés de $\operatorname{comp}(\mathcal A)$ dans la catégorie
des objets spectraux à valeurs dans $D(\mathcal A)$.

\VerdierSemanticAnchor{ch4d-4-2-4-first-canonical}
Soient $X$ un objet de $\operatorname{comp}(\mathcal A)$ et
$X(q)\subset X$, $q\in\mathbb Z$, sa première filtration canonique.
L'objet spectral associé par la proposition (4.2.3) à l'objet filtré
$\bigl(X,X(q),q\in\mathbb Z\bigr)$ n'est autre, dans ce cas, que le premier
objet spectral, à valeurs dans $D(\mathcal A)$, associé au complexe $X$
(4.1.5).

\VerdierSemanticAnchor{ch4d-4-2-4-nfold}
Soient de même $X$ un complexe $n$-uple de $\mathcal A$,
$\int_n X=Y$ le complexe simple associé, $m$ un entier,
$1\leq m\leq n$, et :
\[
0={}^{m}\!Y(-\infty)\subset\cdots\subset{}^{m}\!Y(q)
\subset\cdots\subset{}^{m}\!Y(+\infty)=Y
\]
la filtration définie par le $m$-ième degré de $X$ (4.1.6.9). L'objet
spectral associé par la proposition (4.2.3) à l'objet filtré
$\bigl(Y,{}^{m}\!Y(q),q\in\mathbb Z\bigr)$ n'est autre, dans ce cas, que
l'objet spectral, à valeurs dans $D(\mathcal A)$, défini par $X$, filtré
par son $m$-ième degré (4.1.8).

\VerdierSemanticAnchor{ch4d-4-3}
\medskip\noindent\textbf{4.3. Deuxième objet spectral canonique.}

\VerdierSemanticAnchor{ch4d-4-3-1}
\medskip\noindent\textbf{4.3.1.}\quad
Soit $X$ un objet de $\operatorname{comp}(\mathcal A)$. Posons :

\VerdierSemanticAnchor{ch4d-4-3-1-1}
\begin{equation}
\begin{aligned}
X(-\infty)&=0,\qquad X(+\infty)=X,\\
X(q)&=\cdots\longrightarrow X^{q-3}\longrightarrow X^{q-2}
\longrightarrow\operatorname{Ker}(d_X^{q-1})
\longrightarrow0\longrightarrow0\longrightarrow\cdots,
\qquad q\in\mathbb Z,
\end{aligned}
\tag{4.3.1.1}
\end{equation}
les différentielles de ce dernier complexe étant induites par les
différentielles de $X$. On a des monomorphismes canoniques :

\VerdierSemanticAnchor{ch4d-4-3-1-2}
\begin{equation}
0=X(-\infty)\subset\cdots\subset X(q)\subset X(q+1)
\subset\cdots\subset X(+\infty)=X
\tag{4.3.1.2}
\end{equation}
qui définissent une filtration de type $\widetilde{\mathbb Z}$ sur $X$
qu'on appelle la \emph{deuxième filtration canonique de $X$}. On en déduit,
par la proposition (4.2.3), un objet spectral
$\bigl(X(p,q),\delta(p,q,r)\bigr)$ qu'on appelle le \emph{deuxième objet
spectral, à valeurs dans $D(\mathcal A)$, associé au complexe $X$}.

\VerdierSemanticAnchor{ch4d-4-3-1-associated-functor}
Un morphisme dans $\operatorname{comp}(\mathcal A)$,
$X\longrightarrow Y$, respecte les deuxièmes filtrations canoniques et,
par suite, le foncteur :
\[
X\longmapsto\bigl(X(p,q),\delta(p,q,r)\bigr)
\]
est un foncteur de $\operatorname{comp}(\mathcal A)$ dans la catégorie des
objets spectraux à valeurs dans $D(\mathcal A)$. Ce foncteur est appelé le
\emph{foncteur deuxième objet spectral canonique}.
\endgroup
